\scanpage{24}
\mspage{21}
% archive sheet number: 23
Autrement dit, on peut considérer
\[
\underset{\substack{\rotatebox{-90}{$\scriptstyle\in$}\\\mathrm{Diag}}}{I}
\longmapsto
\underset{\substack{\rotatebox{-90}{$\scriptstyle\in$}\\\mathrm{DER}}}{\mathbb D_I}
\]
comme fonctoriel \deleted{\unsure{en}} en $I$ de trois façons \inserted{(\unsure{une contr.}, les deux autres ``co'')}, suivant qu'on utilise les $f^*$, les $f_!$ (s'ils existent) ou les $f_*$ (s'ils existent). On désignera par
\[
f^{\mathbb D}_!,\quad f^*_{\mathbb D},\quad f^{\mathbb D}_*
\]
indifféremment les foncteurs entre $\mathbb D(I)$ et $\mathbb D(I')$, ou (de façon plus fine) entre $\mathbb D_I$ et $\mathbb D_{I'}$, dont les précédents se déduisent en passant aux catégories fondamentales correspondantes. De $\unsure{=|\simeq}$ pour $f^!$ ou pour $f^?$ sous réserve d'existence. P. ex. si on suppose que les $f^!$ existent pour $\mathbb D$, quel que soit $f$ imm. fermée, alors on trouve \inserted{pour \deleted{\unsure{tout}} tel $f$} un morphisme de dérivateurs
\[
\mathbb D_I\xleftarrow{f^!_{\mathbb D}}\mathbb D_{I'}.
\]

Il s'impose de regarder si tous \unsure{les|ces} morphismes de dérivateurs sont exacts, i.e. si pour
\[
g:J\longrightarrow J'
\]
dans $\mathrm{Diag}$, ils commutent à $g_!$, $\unsure{g_*|J_*}$.
